\documentclass[margin,line]{res}

\usepackage{hyperref}


\oddsidemargin -.5in
\evensidemargin -.5in
\textwidth=6.0in
\itemsep=0in
\parsep=0in
% if using pdflatex:
%\setlength{\pdfpagewidth}{\paperwidth}
%\setlength{\pdfpageheight}{\paperheight} 

\newenvironment{list1}{
  \begin{list}{\ding{113}}{%
      \setlength{\itemsep}{0in}
      \setlength{\parsep}{0in} \setlength{\parskip}{0in}
      \setlength{\topsep}{0in} \setlength{\partopsep}{0in} 
      \setlength{\leftmargin}{0.17in}}}{\end{list}}
\newenvironment{list2}{
  \begin{list}{$\bullet$}{%
      \setlength{\itemsep}{0in}
      \setlength{\parsep}{0in} \setlength{\parskip}{0in}
      \setlength{\topsep}{0in} \setlength{\partopsep}{0in} 
      \setlength{\leftmargin}{0.2in}}}{\end{list}}


\begin{document}

\name{Jamie Hayes \vspace*{.1in}}

\begin{resume}
\section{\sc Contact Information}
\vspace{.05in}
\begin{tabular}{@{}p{3in}p{4in}}
{\it Email:}    \url{jamiedhayes1@gmail.com} 
& {\it Website:} \href{https://scholar.google.com/citations?hl=en&user=95gIFIgAAAAJ}{My research scholar page}\\      
   
\end{tabular}

\section{\sc Experience}

{\bf Google DeepMind}, London, UK

\vspace{-.3cm}
{\em Staff Research Scientist} \hfill {\bf May, 2025 - Present}
\vspace{1em}
\begin{itemize}
\item Research lead for Gemini Security Modelling, a team of 9 members of technical staff working on evaluating and improving Gemini's resilience to prompt injection and jailbreak attacks in security-sensitive settings.
\end{itemize}

{\bf Google DeepMind}, London, UK

\vspace{-.3cm}
{\em Senior Research Scientist} \hfill {\bf May, 2023 - May, 2025}
\vspace{1em}
\begin{itemize}
\item Proposed and led watermarking project, which then became SynthID: \url{https://deepmind.google/technologies/synthid/}. This is Google's provenance detection system that is deployed across all generative models available to Google customers.
\item AI agents that have the ability to act upon untrusted data (for example, a model that can read and write emails) are vulnerable to ``indirect prompt injection'', an attack where the model may share private user data with an untrusted party. I lead a team developing internal benchmarks and auditing techniques to measure the vulnerability to indirect prompt injection attacks, and more generally our remit is to design Trustworthy \& Capable AI Agents. Our work will be deployed across all models developed by Google, reducing the risk of leaking private user data.
\end{itemize}

{\bf Google DeepMind}, London, UK

\vspace{-.3cm}
{\em Research Scientist} \hfill {\bf June, 2020 -  May, 2023}

Research into verification and robustness guarantees in machine learning. Developed techniques to audit privacy-preserving ML pipelines -- now deployed in Google Photos.

{\bf Google DeepMind}, London, UK

\vspace{-.3cm}
{\em Internship} \hfill {\bf May, 2019 - September, 2019}

Research into watermarking applications in machine learning. This eventually led to the formation of the SynthID project.

{\bf Google Research}, Mountain View, CA, USA

\vspace{-.3cm}
{\em Internship} \hfill {\bf July, 2018 - September, 2018}

Research into generative and adversarial machine learning. Developed new techniques for unsupervised style transfer,
and new probabilistic conditioning methods for generative models.

{\bf Microsoft Research}, Cambridge, UK

\vspace{-.3cm}
{\em Internship} \hfill {\bf February, 2018 - May, 2018}

Research into security and privacy attacks and defenses in multi-party machine learning.

{\bf Naval Research Laboratory}, Washington, DC, USA

\vspace{-.3cm}
{\em Internship} \hfill {\bf August, 2017 - October, 2017}

Research into adversarial machine learning and network traffic analysis. Developed a project using machine learning
 to perform privacy-preserving experiments on live traffic analysis on Tor.

{\bf Government Digital Services}, London, UK

\vspace{-.3cm}
{\em Intern for the SecOps team} \hfill {\bf March,
  2017 - July, 2017}

I developed and implemented a privacy-preserving machine learning pipeline to aid threat analysis
and improve Transaction Monitoring (TxM) on the GOV.UK Verify system.

{\bf University of Manchester}, Manchester, UK

\vspace{-.3cm}
{\em Network Technician} \hfill {\bf September, 2011 - March, 2013}

Maintenance of the internal University network that provided a
connection to over 10,000 students. Duties included server
maintenance, switch configurations, some SDN programming.




\section{\sc Education}
{\bf Dept. of Computer Science, University College London}, UK\\
{\bf Sep. 2014 - June. 2020} - Ph.D. candidate\\
Privacy, Security and Machine Learning\\
Advisor:  Prof. George Danezis\\
Second Advisor: Prof. Thore Graepel

{\bf Computer Laboratory, University of Cambridge}, UK\\
{\bf Sep. 2013 - Apr. 2014} - Researcher\\
Parameterized Computational Complexity of the Graph Isomorphism Problem\\
Advisor: Prof. Anuj Dawar

{\bf Dept. of Mathematics, University of Manchester}, UK\\
%{\em Department of Mathematics and Statistics} 
{\bf Sep. 2007 - Jun. 2011} - Master of Mathematics\\
First Class Honours\\
Advisor: Prof. Richard Sharp


% \section{\sc Honors and Awards} 


% %Alan Turing Institute Summer Research Grant, 2016 (Refused)

% %\vspace*{-2.5mm}

% NeurIPS 2018 Student Travel Award\newline

% \vspace*{-3mm}

% NeurIPS 2017 Student Travel Award\newline

% \vspace*{-3mm}

% Google Phd Fellowship in Machine Learning (2017-2020)\newline

% \vspace*{-3mm}

% \noindent Invited to Google PhD Summit in Security (2016) and Machine
% Learning (2017)\newline

% \vspace*{-3mm}

% Academic Center of Excellence Studentship, 2014\newline

% \vspace*{-3mm}

% Engineering and Physical Sciences Research Council (EPSRC) Doctoral
% Training Studentship, 2013


% \section{\sc Research}

% Extensions and limitations of randomized smoothing for robustness guarantees\\
% J Hayes 06-2020 CVPR (workshop track)  

% A framework for robustness certification of smoothed classifiers using f-divergences\\
% K Dvijotham, J Hayes, B Balle, Z Kolter, C Qin, A Gyorgy, K Xiao, S Gowal, P Kohli 05-2020 ICLR

% Towards transformation-resilient provenance detection of digital media\\
% J Hayes, K Dvijotham, Y Chen, S Dieleman, P Kohli, N Casagrande 09-2019  

% LOGAN: Membership inference attacks against generative models\\
% J Hayes, L Melis, G Danezis, E De Cristofaro 07-2019 PETS  

% A note on hyperparameters in black-box adversarial examples\\
% J Hayes 12-2018  

% Evading classifiers in discrete domains with provable optimality guarantees\\
% B Kulynych, J Hayes, N Samarin, C Troncoso 12-2018 NeurIPS (workshop track)  

% Contamination attacks in multi-party machine learning\\
% J Hayes, O Ohrimenko 12-2018 NeurIPS 

% On visible adversarial perturbations \& digital watermarking\\
% J Hayes 06-2018 CVPR (workshop track) 

% Learning universal adversarial perturbations with generative models\\
% J Hayes 05-2018 DLS (IEEE S\&P workshop)  

% Generating steganographic images via adversarial training\\
% J Hayes, G Danezis 12-2017 NeurIPS  

% AnNotify: A private notification service\\
% A Piotrowska, J Hayes, N Gelernter, G Danezis, A Herzberg 10-2017 WPES  

% The Loopix anonymity system\\
% A Piotrowska, J Hayes, T Elahi, S Meiser, G Danezis 08-2017 USENIX Security  

% Website fingerprinting defenses at the application layer\\
% G Cherubin, J Hayes, M Juarez 07-2017 PETS  

% TASP: Towards anonymity sets that persist\\
% J Hayes, C Troncoso, G Danezis 10-2016 WPES  

% \emph{k}-fingerprinting: a robust scalable website fingerprinting technique\\
% J Hayes, G Danezis 08-2016 USENIX Security  

% Traffic confirmation attacks despite noise\\
% J Hayes 02-2016 NDSS (workshop track)  

% Guard sets for onion routing\\
% J Hayes 07-2015 PETS  

% An introduction to the dynamics of real and complex quadratic polynomials\\
% J Hayes 07-2011  

% \section{\sc Presentations}

% On visible adversarial perturbations \& digital watermarking,
% CVPR Workshop track

% Learning universal adversarial perturbations with generative models,
% DLS (IEEE S\&P Workshop)

% Invited talk on Adversarial Machine Learning\\
% IBM, IBM Thomas J. Watson Research Center

% Invited talk on Network Traffic Analysis\\
% UK Gov

% TASP: Towards Anonymity Sets that Persist\\
% WPES (Workshop of CCS) 

% \emph{k}-fingerprinting: a Robust Scalable Website Fingerprinting
% Technique\\
% USENIX Security 

% Traffic Confirmation Attacks Despite Noise\\
% Understanding and Enhancing Online Privacy (Workshop of NDSS)

% Guard Sets for Onion Routing\\
% PETS 

% Guard Sets for Onion Routing\\
% University College London Information Security Seminar

% Secure Sets in Graphs\\
% Programming, Logic, and Semantics Reading Group, Computer Lab, University of Cambridge


% \section{\sc Program committee member}

% AAAI 2020

% CVPR 2020

% NeurIPS 2019, 2020

% ICLR 2020 Workshop on "Towards Trustworthy ML: Rethinking Security and Privacy for ML"

% PrivateNLP'20 (WSDM 2020 Workshop on Privacy and Natural Language Processing)

% ICML 2019

%  PPML19 (Privacy-Preserving Machine Learning - CCS 2019 Workshop)

% Privacy Enhancing Technologies Symposium 2018, 2019, 2020, 2021

% NeurIPS 2018 Workshop on Security in Machine Learning

% (External) CCS, 2017

% (External) NDSS, 2016

% (External) IEEE Symposium on Security and Privacy, 2016

% Transactions on Information Forensics \& Security

% Transactions on Pattern Analysis and Machine Intelligence

% \section{\sc Teaching}

% Advisor (along with George Danezis) for Hugo Picq's MSc Thesis, `Understanding and Scaling a Robust Neural Network', 2019

% Advisor for Axel Goetz's MSc Thesis, `Evaluating the use of Deep Learning for Website Fingerprinting', 2017


\section{\sc Skills} 
\textbf{Knowledge of}:\\
Python \textbullet{} Go \textbullet{} Shell \textbullet{} \LaTeX \textbullet{}
Unix/Linux \textbullet{} Vim \textbullet{} Machine Learning 
(TensorFlow, PyTorch, JAX, SKLearn) \textbullet{} C++ \textbullet{}   Git

\textbf{Exposure to}:\\
Haskell \textbullet{} JavaScript \textbullet{} C \textbullet{} SQL \textbullet{} HTML
\textbullet{} CSS 

\section{\sc Code and Side Projects}

Public code available at \url{https://github.com/jhayes14}, code for
private projects available on request. 

\section{\sc References} 

Available upon request.

\section{\sc Services} 

\begin{itemize}
    \item Reviewing for most major academic ML and Security conferences since 2015 (IEEE S\&P, CCS, USENIX Security, PETS, ICML, ICLR, NeurIPS etc.). Full breakdown available upon request. 
    \item Organizer of the workshop Synthetic Data Generation: Quality, Privacy, Bias at ICLR 2021.
    \item Various invited talks and lectures at conferences, academic labs, summer schools etc. Full breakdown available upon request. 
    
\end{itemize}

% \textbf{George Danezis}\newline
% Reader in Security and Privacy Engineering (Professor) \newline
% Email: g.danezis[at]ucl.ac.uk \newline
% University College London,\newline
% Dept. of Computer Science,\newline
% Gower Street, London WC1E 6BT, U.K. \newline

% \textbf{Louise Walker}\newline
% Reader in Mathematics \newline
% Email: Louise.Walker[at]manchester.ac.uk \newline
% Room 2.243, Alan Turing Building \newline
% School of Mathematics,\newline
% University of Manchester \newline
% Oxford Road, Manchester M13 9PL, UK \newline


\end{resume}
\end{document}



